\chapter{Topology, poles and the Coxeter presentation}
\label{ch:ArchonFull600Theory_Topology}

This chapter is Milestone 4 of the handoff: Lemma 5 (the inradius and the
cell-pole normalization), Lemma 2 (the radial topology of the boundary,
with the standard topology E3 as a hypothesis) and the Coxeter
presentation of Lemma 4 (with E5 as a hypothesis, in the form of the
order of \(W(H_4)\)). Lemma 5 comes first because Lemma 2 uses its inscribed
ball. Source line numbers refer to
\texttt{research/Full\_600cell\_Puzzle\_Theory.tex}.

\section{Facet poles and the inradius}

\begin{lemma}[Facet poles]
  \label{lem:facet-pole}
  \lean{ArchonFull600Theory.facetPole, ArchonFull600Theory.inradius,
    ArchonFull600Theory.norm_sq_facetPole, ArchonFull600Theory.norm_facetPole,
    ArchonFull600Theory.facetPlane_iff}
  \leanok
  \uses{lem:supporting-facet}
  % SOURCE: Lemma 5, lines 708--718
  % SOURCE QUOTE: \|N\|^2=4(4+6\phi), \end{equation} the perpendicular foot of the origin on
  % its plane is \begin{equation} n=\frac{N}{4}. [...] \|n\|^2=\frac{4+6\phi}{4}
  % =\frac{\phi^4}{2}.
  \textit{Source: Lemma 5.}
  For a four-clique with vertex sum \(N\), the pole \(n=N/4\) has
  \(\|n\|^2=(4+6\phi)/4=\phi^4/2\), so \(\|n\|=\phi^2/\sqrt2\), and the
  facet plane \(N\cdot x=4+6\phi\) is the plane \(n\cdot x=\|n\|^2\).
\end{lemma}
\begin{proof}
  \leanok
  \(\|N\|^2=16+24\phi\) is the supporting-facet computation, and
  \(\phi^4=3\phi+2\).
\end{proof}

\begin{theorem}[\(P\) is cut out by its facets]
  \label{thm:half-spaces}
  \lean{ArchonFull600Theory.exists_flip, ArchonFull600Theory.exists_facetCone,
    ArchonFull600Theory.polytope_eq_halfSpaces, ArchonFull600Theory.polytope_eq_poleHalfSpaces}
  \leanok
  \uses{lem:supporting-facet, thm:face-counts}
  \textit{Not in the source as a separate statement; it supports the
  inradius.}
  \(P\) is the intersection of the half-spaces \(N\cdot x\le4+6\phi\) of
  the \(600\) four-cliques, equivalently of the half-spaces
  \(n\cdot x\le\|n\|^2\) of their poles.
\end{theorem}
\begin{proof}
  \leanok
  \(P\) lies in each half-space by the supporting-facet test. Conversely,
  let \(x\) satisfy every inequality and take a four-clique whose normal
  maximizes \(N\cdot x\). Its four vertices are a basis of \(\mathbb H\).
  If the coefficient of a vertex \(v_j\) were negative, the triangle
  opposite \(v_j\) has a second common neighbour \(v_{j'}\), and the
  four-clique across that triangle has \(N'-N=v_{j'}-v_j\), which is
  orthogonal to the three shared vertices. Then
  \(N'\cdot x-N\cdot x=c_j(v_{j'}\cdot v_j-4)>0\), against maximality. So
  \(x\) is a nonnegative combination of the four vertices, and
  \(N\cdot x\le4+6\phi\) bounds the sum of the coefficients by \(1\). Hence
  \(x\) lies in the convex hull of the four vertices and \(0\).
\end{proof}

\begin{theorem}[Inradius]
  \label{thm:inradius}
  \lean{ArchonFull600Theory.closedBall_inradius_subset, ArchonFull600Theory.isGreatest_inradius}
  \leanok
  \uses{lem:facet-pole, thm:half-spaces}
  % SOURCE: Lemma 5, lines 695--697 and 719
  % SOURCE QUOTE: The radius-two model has inradius \begin{equation}
  % r_{\mathrm{in}}=\frac{\phi^2}{\sqrt2}. [...] Regularity gives the same inradius at every
  % facet.
  \textit{Source: Lemma 5.}
  \(\phi^2/\sqrt2\) is the largest radius of a closed ball about \(0\)
  inside \(P\).
\end{theorem}
\begin{proof}
  \leanok
  Every pole has norm \(\phi^2/\sqrt2\), so the Cauchy--Schwarz inequality
  puts the ball in every pole half-space, hence in \(P\). A larger ball
  contains a point beyond a pole on its ray, which violates that pole's
  inequality. The source appeals to regularity; here every pole norm is
  computed.
\end{proof}

\begin{theorem}[Lemma 5: cell-pole normalization]
  \label{thm:cell-pole}
  \lean{ArchonFull600Theory.cellPole, ArchonFull600Theory.cellPoleMap,
    ArchonFull600Theory.exists_rotation_eq, ArchonFull600Theory.cellPoleNormalization}
  \leanok
  \uses{lem:facet-pole, thm:half-spaces, thm:inradius, def:lpq, lem:lpq-det}
  % SOURCE: Lemma 5, lines 693--735
  % SOURCE QUOTE: A similarity with scale factor \(4/\phi\) and a suitable rotation places a
  % chosen facet pole at \begin{equation} n_0=(\phi^3,1,1,1), \end{equation} with every
  % facet plane written as \begin{equation} n_c\cdot x=\|n_c\|^2.
  \textit{Source: Lemma 5.}
  The cell pole \(n_0=(\phi^3,1,1,1)\) has
  \(\|n_0\|^2=\phi^6+3=8\phi^2\) and \(\|n_0\|/r_{\mathrm{in}}=4/\phi\). For
  each chosen four-clique there is \(Q\in SO(4)\) such that
  \(\sigma x=(4/\phi)Qx\) sends its pole to \(n_0\), every pole \(n\) to a
  pole \(\tilde n=\sigma n\) of norm \(\|n_0\|\), every facet plane
  \(n\cdot x=\|n\|^2\) to the plane \(\tilde n\cdot y=\|\tilde n\|^2\) and
  every facet to the face of \(\sigma P\) on that plane; and \(\sigma P\)
  is the intersection of the half-spaces
  \(\tilde n\cdot y\le\|\tilde n\|^2\). The remark that the symmetry group
  of \(\sigma P\) is the conjugate of \(K\) is not stated.
\end{theorem}
\begin{proof}
  \leanok
  \(\phi^6=8\phi+5\) gives \(\phi^6+3=8\phi^2\), and the ratio follows.
  For nonzero \(x,y\), the map \(L_{b,a}\) with \(a=x/\|x\|\) and
  \(b=y/\|y\|\) lies in \(SO(4)\) and sends \(x\) to \((\|x\|/\|y\|)y\);
  with \(x\) the chosen pole and \(y=n_0\), the scale \(4/\phi\) sends the
  pole to \(n_0\). A similarity multiplies inner products by
  \((4/\phi)^2\), so it preserves the plane and half-space equations.
\end{proof}

\section{The radial topology of the boundary}

\begin{theorem}[Radial homeomorphism]
  \label{thm:radial-homeomorph}
  \lean{ArchonFull600Theory.polytopeRescale, ArchonFull600Theory.boundaryHomeomorph,
    ArchonFull600Theory.boundaryHomeomorph_apply}
  \leanok
  \uses{thm:inradius}
  % SOURCE: Lemma 2, lines 402--423
  % SOURCE QUOTE: The map \begin{equation} \partial P\longrightarrow S^3,\qquad
  % x\longmapsto\frac{x}{\|x\|} \end{equation} is therefore a continuous bijection.
  \textit{Source: Lemma 2.}
  \(0\) is an interior point of \(P\), \(P\) is compact and lies in the
  ball of radius \(2\), and \(x\mapsto x/\|x\|\) is a homeomorphism from
  \(\partial P\) onto the unit sphere \(S^3\subseteq\mathbb H\).
\end{theorem}
\begin{proof}
  \leanok
  The inscribed ball of Lemma 5 makes \(0\) interior; the source uses the
  cross-polytope \(\operatorname{conv}\{\pm2e_i\}\). Mathlib's gauge
  rescaling homeomorphism of the convex body \(P\) to the closed unit ball
  maps the frontier onto the unit sphere and, since the gauge of \(P\) is
  \(1\) on its frontier, agrees there with \(x\mapsto x/\|x\|\). The
  homeomorphism is constructed, so the compact-to-Hausdorff step is not
  needed.
\end{proof}

\begin{definition}[E3 as a hypothesis]
  \label{def:sphere-topology}
  \lean{ArchonFull600Theory.intHomology, ArchonFull600Theory.SphereTopology}
  \leanok
  % SOURCE: Lemma 2, lines 425--430
  % SOURCE QUOTE: The sphere \(S^3\) has a CW decomposition with one zero-cell and one
  % three-cell, giving the stated integral homology. [...] The van Kampen theorem gives
  % trivial fundamental group.
  \textit{Source: E3 in the handoff's table of external results.}
  \(S^3\) is simply connected, and its integral singular homology is
  \(\mathbb Z\) in degrees \(0\) and \(3\) and zero in every other degree.
\end{definition}

\begin{theorem}[Lemma 2]
  \label{thm:radial-topology}
  \lean{ArchonFull600Theory.intHomology_iso, ArchonFull600Theory.radialTopology}
  \leanok
  \uses{thm:radial-homeomorph, def:sphere-topology}
  % SOURCE: Lemma 2, lines 388--398
  % SOURCE QUOTE: The boundary \(\partial P\) is homeomorphic to \(S^3\).
  \textit{Source: Lemma 2.}
  Under E3, \(H_j(\partial P;\mathbb Z)\) is \(\mathbb Z\) for \(j=0,3\)
  and zero otherwise, and \(\pi_1(\partial P)=1\).
\end{theorem}
\begin{proof}
  \leanok
  Singular homology and simple connectivity are invariant under the
  homeomorphism of Theorem~\ref{thm:radial-homeomorph}.
\end{proof}

\section{The Coxeter presentation}

\begin{definition}[Base flag and wall reflections]
  \label{def:walls}
  \lean{ArchonFull600Theory.flagFace, ArchonFull600Theory.flagBarycenter,
    ArchonFull600Theory.wallVec, ArchonFull600Theory.wallRefl, ArchonFull600Theory.wallPerm}
  \leanok
  % SOURCE: Lemma 4, lines 586--596
  % SOURCE QUOTE: A flag gives a spherical tetrahedral chamber. [...] For each of the four
  % chamber walls, the symmetry taking the chosen chamber to its adjacent chamber fixes the
  % other three barycenter rays. It is the reflection in their span.
  \textit{Source: Lemma 4, proof of the presentation.}
  The base flag is
  \(\{v_0\}\subset\{v_0,v_{30}\}\subset\{v_0,v_{30},v_{31}\}\subset
  \{v_0,v_{30},v_{31},v_{32}\}\). For rank \(i\), \(s_i\) is the
  reflection in the hyperplane orthogonal to \(v_a-v_b\), with
  \((a,b)=(0,30),(30,31),(31,32),(32,42)\); vertex \(42\) is the second
  common neighbour of the triangle \(\{0,30,31\}\).
\end{definition}

\begin{theorem}[Chambers and walls]
  \label{thm:chamber-walls}
  \lean{ArchonFull600Theory.wallRefl_vertexR, ArchonFull600Theory.wallRefl_mem,
    ArchonFull600Theory.span_flagBarycenter, ArchonFull600Theory.eq_one_of_fix_flagBarycenter,
    ArchonFull600Theory.span_flagBarycenter_compl, ArchonFull600Theory.chamberWalls}
  \leanok
  \uses{def:walls, def:symmetry-group, thm:face-counts}
  % SOURCE: Lemma 4, lines 586--596
  % SOURCE QUOTE: A symmetry fixing a flag fixes the barycenters of its four faces; these
  % four vectors are linearly independent, so the symmetry is the identity.
  \textit{Source: Lemma 4, proof of the presentation.}
  The flag is a chain of cliques, its four barycenters span \(\mathbb H\),
  and a linear isometry fixing them is the identity. Each \(s_i\) lies in
  \(K\), fixes the faces of the other ranks and moves the face of rank
  \(i\), and is the reflection in the span of the other three
  barycenters.
\end{theorem}
\begin{proof}
  \leanok
  For \(u=v_a-v_b\), the vector \((u\cdot u)v-2(v\cdot u)u\) is computed
  exactly in \(\mathbb Z[\phi]\) for each vertex \(v\) and found to be
  \((u\cdot u)\) times a vertex; four kernel-checked tables record the
  resulting permutations. The face images and the fixed barycenters follow
  from the tables, and a dimension count identifies the mirror with the
  span of the other three barycenters.
\end{proof}

\begin{theorem}[Relations and generation]
  \label{thm:coxeter-relations}
  \lean{ArchonFull600Theory.H₄_chain, ArchonFull600Theory.wallRefl_mul_pow,
    ArchonFull600Theory.orderOf_wallRefl_mul, ArchonFull600Theory.closure_wallRefl}
  \leanok
  \uses{thm:chamber-walls, thm:finite-symmetry-group, thm:binary-icosahedral}
  % SOURCE: Lemma 4, lines 598--601
  % SOURCE QUOTE: Chamber adjacency is connected, so these reflections generate \(K\).
  % Consecutive rank changes give rank-two sections of sizes \(3,3,5\), producing the
  % displayed relations. Nonconsecutive changes commute and give square galleries.
  \textit{Source: Lemma 4, proof of the presentation.}
  Mathlib's \(H_4\) Coxeter matrix is the chain \((3,3,5)\);
  \((s_is_j)^{m_{ij}}=1\) and \(s_is_j\) has order exactly \(m_{ij}\); and
  \(s_0,\ldots,s_3\) generate \(K\).
\end{theorem}
\begin{proof}
  \leanok
  An element of \(K\) is determined by the images of the four vertices of
  the base facet, so the relations and the minimality of each order are
  kernel checks on four vertices. For generation, explicit words in the
  \(s_i\), found by a breadth-first search outside Lean and checked by the
  kernel, equal \(L_{v_8/2,1}\), \(L_{v_{30}/2,1}\), \(L_{1,v_8/2}\),
  \(L_{1,v_{30}/2}\) and \(J\). The word tree of Milestone 3 writes every
  element of \(I^*\) in \(v_8/2\) and \(v_{30}/2\), and every element of
  \(K\) is \(L_{p,q}\) or \(L_{p,q}J\). The source argues through the
  connectivity of chamber adjacency instead.
\end{proof}

\begin{definition}[E5 as a hypothesis]
  \label{def:coxeter-sources}
  \lean{ArchonFull600Theory.CoxeterSources}
  \leanok
  % SOURCE: Lemma 4, lines 603--612
  % SOURCE QUOTE: Hence the displayed presentation presents \(K\), rather than merely mapping
  % onto it. The Coxeter group with this chain presentation is \(W(H_4)\).
  \textit{Source: E5 in the handoff's table of external results, in another
  form.}
  The Coxeter group presented by the chain \((3,3,5)\) has order
  \(14400\). The source instead proves completeness of the presentation by
  closed galleries in the dual cell complex of \(S^3\) (E5's second part,
  with Lemma 2); that argument is not formalized.
\end{definition}

\begin{theorem}[Lemma 4: the Coxeter presentation]
  \label{thm:coxeter-presentation}
  \lean{ArchonFull600Theory.coxeterHom, ArchonFull600Theory.coxeterHom_surjective,
    ArchonFull600Theory.coxeterHom_bijective, ArchonFull600Theory.coxeterPresentation}
  \leanok
  \uses{thm:coxeter-relations, def:coxeter-sources, thm:finite-symmetry-group}
  % SOURCE: Lemma 4, lines 533--543
  % SOURCE QUOTE: Moreover, \begin{equation} K\cong W(H_4) \end{equation} with presentation
  \textit{Source: Lemma 4.}
  The Coxeter group of type \(H_4\) maps onto \(K\) with \(s_i\) going to
  the wall reflections. Under E5, in the form of
  Definition~\ref{def:coxeter-sources}, this map is an isomorphism
  \(W(H_4)\cong K\).
\end{theorem}
\begin{proof}
  \leanok
  The relations give the homomorphism through the universal property of
  the Coxeter group, and generation makes it onto. Both groups have
  \(14400\) elements, so the surjection is a bijection.
\end{proof}
